\section{AI 工具使用声明（提案）}

本队在论文结构梳理、文字校对、证据链检查和可复现性清单整理中使用了 AI 辅助工具；核心问题理解、变量定义、模型设计、程序实现、数据处理、求解运行、结果核验和最终取舍由参赛队独立完成。AI 输出不作为数据来源、算法运行结果或数学结论的替代品，所有进入论文的数字均回溯至冻结主张及其哈希锁定的实验文件。

在正式提交前，根 Agent 应依据 `state/decision_log.json` 中的 `compliance.ai_usage` 自动生成并核对支持材料 `AI工具使用详情.pdf`，逐项列出工具名称、提供方、模型或版本、使用日期、使用阶段、关键提示词与响应定位、实际采纳内容以及人工修改和验证记录。未经 usage ledger 核验的字段不得手工补写。论文中由 AI 辅助形成或改写的段落应在对应位置按赛事要求标注；参考文献中仅列出实际使用且已核验的工具信息，不把 AI 工具当作学术证据引用。

本提案对应的问题二至问题四仅协助表达已经冻结的模型接口、结果和边界，未新增实验、未改变冻结主张，也未把启发式、滚动 MILP、顺序耦合或离散峰值扫描表述为其证据范围之外的全局最优结论。
